% ------------------------------------------------------------------------------
% Resultados
% ------------------------------------------------------------------------------

\chapter{Resultados}

Cada capítulo deve conter uma pequena introdução (tipicamente, um ou dois parágrafos), que deve deixar claro o objetivo e o que será discutido no capítulo, bem como a organização do mesmo.

Neste capítulo, são apresentados e analisados os resultados obtidos a partir das simulações computacionais do modelo de queda livre. São exibidos os principais gráficos gerados, bem como tabelas que sintetizam os dados numéricos relevantes para a avaliação do comportamento do sistema e da acurácia dos métodos utilizados. A seguir, são discutidos os resultados à luz dos objetivos propostos.




\section{Apresentação dos Gráficos}
\label{sec_graficos_resultados}


Os gráficos a seguir ilustram a evolução das principais variáveis do sistema ao longo do tempo, bem como a comparação entre diferentes cenários simulados.




\section{Resultados para o Modelo de uma Partícula em Queda Livre}
Foram realizados experimentos simulando a queda livre de uma partícula, variando o passo de tempo ($\Delta t$) para analisar o impacto na precisão dos resultados. Os gráficos a seguir apresentam, respectivamente, a posição, velocidade e aceleração em função do tempo para diferentes valores de $\Delta t$.
\begin{figure}[H]
	\centering
	\begin{subfigure}[b]{0.32\textwidth}
		\includegraphics[width=\textwidth]{./figuras/graficos/queda_livre_particula/grafico_posicao_tempo_multidt}
		\caption{Posição}
	\end{subfigure}
	\hfill
	\begin{subfigure}[b]{0.32\textwidth}
		\includegraphics[width=\textwidth]{./figuras/graficos/queda_livre_particula/grafico_velocidade_tempo_multidt}
		\caption{Velocidade}
	\end{subfigure}
	\hfill
	\begin{subfigure}[b]{0.32\textwidth}
		\includegraphics[width=\textwidth]{./figuras/graficos/queda_livre_particula/grafico_aceleracao_tempo_multidt}
		\caption{Aceleração}
	\end{subfigure}
	\caption{Resultados da simulação da queda livre de uma partícula para diferentes passos de tempo: (a) posição, (b) velocidade e (c) aceleração em função do tempo.}
\end{figure}

\section{Resultados para o Modelo Corpo Rígido com Resistência Linear}
Para o corpo rígido com resistência linear (Lei de Stokes), os gráficos mostram a evolução das variáveis para diferentes passos de tempo.
\begin{figure}[H]
	\centering
	\begin{subfigure}[b]{0.32\textwidth}
		\includegraphics[width=\textwidth]{./figuras/graficos/queda_livre_corpo_rigido/linear/grafico_posicao_tempo_linear}
		\caption{Posição}
	\end{subfigure}
	\hfill
	\begin{subfigure}[b]{0.32\textwidth}
		\includegraphics[width=\textwidth]{./figuras/graficos/queda_livre_corpo_rigido/linear/grafico_velocidade_tempo_linear}
		\caption{Velocidade}
	\end{subfigure}
	\hfill
	\begin{subfigure}[b]{0.32\textwidth}
		\includegraphics[width=\textwidth]{./figuras/graficos/queda_livre_corpo_rigido/linear/grafico_aceleracao_tempo_linear}
		\caption{Aceleração}
	\end{subfigure}
	\caption{Resultados da simulação do corpo rígido com resistência linear: (a) posição, (b) velocidade e (c) aceleração em função do tempo.}
\end{figure}

\section{Resultados para o Modelo Corpo Rígido com Resistência Quadrática}
Para o corpo rígido com resistência quadrática, os gráficos a seguir apresentam a posição, velocidade e aceleração em função do tempo para diferentes passos de tempo.
\begin{figure}[H]
	\centering
	\begin{subfigure}[b]{0.32\textwidth}
		\includegraphics[width=\textwidth]{./figuras/graficos/queda_livre_corpo_rigido/quadratica/grafico_posicao_tempo_quadratica_multidt}
		\caption{Posição}
	\end{subfigure}
	\hfill
	\begin{subfigure}[b]{0.32\textwidth}
		\includegraphics[width=\textwidth]{./figuras/graficos/queda_livre_corpo_rigido/quadratica/grafico_velocidade_tempo_quadratica_multidt}
		\caption{Velocidade}
	\end{subfigure}
	\hfill
	\begin{subfigure}[b]{0.32\textwidth}
		\includegraphics[width=\textwidth]{./figuras/graficos/queda_livre_corpo_rigido/quadratica/grafico_aceleracao_tempo_quadratica_multidt}
		\caption{Aceleração}
	\end{subfigure}
	\caption{Resultados da simulação do corpo rígido com resistência quadrática: (a) posição, (b) velocidade e (c) aceleração em função do tempo.}
\end{figure}

\section{Resultados da Acurácia dos Métodos Numéricos}
Por fim, a acurácia dos métodos numéricos foi avaliada para cada modelo, variando o passo de tempo. Os gráficos a seguir mostram o erro relativo em função de $\Delta t$ para cada caso e a comparação geral.
\begin{figure}[H]
    \centering
    \begin{subfigure}[b]{0.32\textwidth}
        \includegraphics[width=\textwidth]{./figuras/graficos/acuracia/acuracia_particula}
        \caption{Partícula sem resistência do ar}
    \end{subfigure}
    \hfill
    \begin{subfigure}[b]{0.32\textwidth}
        \includegraphics[width=\textwidth]{./figuras/graficos/acuracia/acuracia_linear}
        \caption{Corpo rígido com resistência linear}
    \end{subfigure}
    \hfill
    \begin{subfigure}[b]{0.32\textwidth}
        \includegraphics[width=\textwidth]{./figuras/graficos/acuracia/acuracia_quadratica}
        \caption{Corpo rígido com resistência quadrática}
    \end{subfigure}
    \caption{Acurácia do método numérico para (a) partícula sem resistência do ar, (b) corpo rígido com resistência linear e (c) corpo rígido com resistência quadrática.}
\end{figure}

A análise dos resultados obtidos demonstra a eficácia dos métodos numéricos empregados na simulação da queda livre sob diferentes regimes físicos. Observou-se que, para passos de tempo suficientemente pequenos, os métodos apresentaram excelente concordância com as soluções analíticas, validando a abordagem computacional adotada. A inclusão dos efeitos de resistência do ar, tanto linear quanto quadrática, permitiu avaliar o impacto desses fatores no comportamento do sistema, evidenciando a importância de modelos mais realistas para a descrição de fenômenos físicos. Por fim, a análise da acurácia reforçou a necessidade de escolha criteriosa do passo de tempo para garantir precisão e estabilidade nos resultados.

